\section{Mengenali Pola Hasil Disassembly dan Metrik Optimasi}
\label{sec:pola-disassembly}

Hasil \eng{disassembly} pada awalnya tampak seperti deretan instruksi
yang tidak berhubungan. Setelah beberapa kali berlatih, mahasiswa
akan menyadari bahwa kompilator dan programmer menghasilkan pola yang
berulang. Mengenali pola tersebut memungkinkan seseorang membaca
program tanpa kode sumber: struktur perulangan, pemilihan,
pemanggilan prosedur, dan interaksi dengan sistem operasi dapat
ditebak dari bentuk instruksinya.

\subsection{Pola-Pola yang Sering Dijumpai}

Tabel~\ref{tab:pola-disassembly} merangkum sembilan pola yang paling
sering muncul, lengkap dengan bentuk hasil \eng{disassembly} dan
cara mengenalinya. Kolom ketiga sengaja memusatkan perhatian pada
\emph{ciri} yang membedakan pola tersebut dari pola lain, karena
itulah yang dipakai saat membaca program yang belum dikenal.

\begin{table}[H]
    \centering
    \caption{Pola kode sumber dan bentuknya pada hasil
    \eng{disassembly}.}
    \label{tab:pola-disassembly}
    \small
    \begin{tabular}{@{}p{2.7cm}p{4.9cm}p{6.9cm}@{}}
        \toprule
        \textbf{Pola kode sumber} & \textbf{Bentuk hasil
        \eng{disassembly}} & \textbf{Cara mengenalinya} \\
        \midrule
        Pembandingan dan percabangan & \lstasm|CMP AX,0064h| +
        \lstasm|JE 0140h| & \asm{CMP} atau \asm{TEST} selalu diikuti
        lompatan bersyarat \\
        \addlinespace
        Perulangan berpenghitung & \lstasm|LOOP 0109h| atau
        \lstasm|DEC CX| + \lstasm|JNZ| & Lompatan mundur yang
        dikendalikan \reg{CX} \\
        \addlinespace
        Perulangan uji-di-akhir & \lstasm|CMP SI,0014h| +
        \lstasm|JB 0100h| & Lompatan mundur setelah pembandingan
        batas \\
        \addlinespace
        \eng{Prologue} prosedur & \lstasm|PUSH BP| +
        \lstasm|MOV BP,SP| & Menyimpan bingkai \eng{stack} lama lalu
        menyediakan variabel lokal \\
        \addlinespace
        \eng{Epilogue} prosedur & \lstasm|POP BP| + \lstasm|RET| &
        \asm{RET} selalu menjadi instruksi terakhir \\
        \addlinespace
        Inisialisasi segmen & \lstasm|MOV AX,1D2Fh| +
        \lstasm|MOV DS,AX| & Pemindahan ke register segmen selalu
        melalui register umum \\
        \addlinespace
        Pemanggilan interupsi & \lstasm|MOV AH,09h| +
        \lstasm|INT 21h| & Nomor fungsi dimuat ke \reg{AH} tepat
        sebelum \asm{INT} \\
        \addlinespace
        Perkalian pangkat dua & \lstasm|SHL AX,1| atau
        \lstasm|SHL AX,CL| & Pergeseran menggantikan \asm{MUL}; tidak
        ada operand pengali \\
        \addlinespace
        Pengosongan register & \lstasm|XOR AX,AX| & Kedua operand
        sama; hasilnya selalu nol \\
        \bottomrule
    \end{tabular}
\end{table}

Hubungan pola-pola tersebut dengan materi bab sebelumnya dapat
ditelusuri kembali: pembandingan dan percabangan pada
Bab~\ref{chap:kontrol}, inisialisasi segmen pada
Bab~\ref{chap:segmentasi}, serta \eng{prologue} dan \eng{epilogue}
prosedur pada Bab~\ref{chap:prosedur}.

Kode~\ref{lst:pola-bongkar} memperlihatkan potongan hasil
\eng{disassembly} yang memuat tiga pola sekaligus: inisialisasi
segmen, perulangan dengan penghitung, dan pemanggilan interupsi.

\begin{session}[caption={Tiga pola dalam satu potongan hasil
\eng{disassembly}.}, label={lst:pola-bongkar}]
1D2F:0100 B82F1D          MOV     AX,1D2F
1D2F:0103 8ED8            MOV     DS,AX
1D2F:0105 B90A00          MOV     CX,000A
1D2F:0108 0304            ADD     AX,[SI]
1D2F:010A 83C602          ADD     SI,+02
1D2F:010D E2F9            LOOP    0108
1D2F:010F B409            MOV     AH,09
1D2F:0111 BA1001          MOV     DX,0110
1D2F:0114 CD21            INT     21
\end{session}

Perhatikan bahwa perintah \code{LOOP 0108} menunjuk ke alamat
\hex{0108}, yaitu alamat yang lebih kecil daripada dirinya sendiri
(\hex{010D}). Lompatan mundur semacam ini adalah tanda perulangan
yang paling dapat diandalkan. Sebaliknya, lompatan maju yang jauh
biasanya menandai percabangan pemilihan.

\subsection{Metrik Optimasi}

Sebelum memperbaiki sesuatu, ukuran keberhasilan harus ditetapkan.
Empat metrik berikut dipakai sepanjang bab ini.

\begin{enumerate}
    \item \textbf{Jumlah instruksi.} Banyaknya instruksi yang
          \emph{dieksekusi} (bukan yang tertulis di kode). Metrik ini
          diperoleh dengan menghitung iterasi.
    \item \textbf{Perkiraan jumlah siklus.} Jumlah siklus \eng{clock}
          yang dibutuhkan. Nilainya dihitung dari tabel siklus tiap
          instruksi, bukan diukur langsung.
    \item \textbf{Ukuran kode dalam byte.} Panjang kode mesin yang
          dihasilkan. Metrik ini penting pada sistem tertanam dengan
          memori terbatas.
    \item \textbf{Jumlah akses memori.} Banyaknya pembacaan dan
          penulisan ke memori. Akses memori jauh lebih mahal daripada
          akses register, sehingga pengurangan satu akses saja sudah
          berarti.
\end{enumerate}

Tabel~\ref{tab:siklus-8086-perkiraan} memuat perkiraan siklus
beberapa instruksi 8086 yang paling sering muncul. Angka-angka ini
berasal dari lembar data Intel dan berlaku untuk kondisi ideal: tanpa
\eng{wait state}, operand di register, dan \eng{prefetch queue} tidak
pernah kosong.

\begin{table}[H]
    \centering
    \caption{Perkiraan siklus instruksi 8086 yang sering dipakai.}
    \label{tab:siklus-8086-perkiraan}
    \small
    \begin{tabular}{@{}p{4.6cm}p{1.7cm}p{2.4cm}p{4.2cm}@{}}
        \toprule
        \textbf{Instruksi} & \textbf{Byte} &
        \textbf{Siklus (perkiraan)} & \textbf{Keterangan} \\
        \midrule
        \lstasm|MOV reg, reg| & 2 & 2 & Paling murah \\
        \addlinespace
        \lstasm|MOV reg, imm16| & 3 & 4 & Memuat konstanta \\
        \addlinespace
        \lstasm|ADD reg, reg| & 2 & 3 & Aritmatika register \\
        \addlinespace
        \lstasm|XOR reg, reg| & 2 & 3 & Sama dengan \asm{MOV reg,0}
        tetapi lebih pendek \\
        \addlinespace
        \lstasm|ADD AX, [SI]| & 2 & 9 + 5 = 14 & EA \lstasm|[SI]| = 5
        siklus \\
        \addlinespace
        \lstasm|ADD AX, [SI+2]| & 3 & 9 + 8 = 17 & EA dengan
        \eng{displacement} lebih mahal \\
        \addlinespace
        \lstasm|MOV mem, reg| & 4 & 9 + 6 = 15 & Alamat langsung \\
        \addlinespace
        \lstasm|SHL reg, 1| & 2 & 2 & Pergeseran satu bit \\
        \addlinespace
        \lstasm|SHL reg, CL| & 2 & 8 + 4$\cdot$CL & Naik seiring
        nilai \reg{CL} \\
        \addlinespace
        \lstasm|MUL reg16| & 2 & 118--133 & Sangat mahal \\
        \addlinespace
        \lstasm|DIV reg16| & 2 & $\approx$150 & Paling mahal \\
        \addlinespace
        \lstasm|LOOP label| & 2 & 17 (lompat) / 5 (berhenti) &
        Overhead perulangan \\
        \addlinespace
        \lstasm|PUSH reg| & 1 & 11 & Akses \eng{stack} \\
        \addlinespace
        \lstasm|LEA reg, [SI+2]| & 3 & 2 + 8 = 10 & Tidak mengubah
        \eng{flag}; lebih ringkas daripada \asm{MOV} + \asm{ADD} \\
        \bottomrule
    \end{tabular}
\end{table}

\begin{infobox}
Angka pada Tabel~\ref{tab:siklus-8086-perkiraan} hanya berlaku bila
operand memang berada di register atau memori seperti tertulis.
Penambahan \eng{wait state}, kegagalan \eng{prefetch}, dan akses ke
memori yang lambat dapat menggandakan waktu sesungguhnya. Itulah
sebabnya satu-satunya bukti optimasi yang sah adalah pengukuran pada
mesin yang dipakai, bukan perhitungan di atas kertas. Hal ini sejalan
dengan prosedur pengukuran pada Bab~\ref{chap:debugging}.
\end{infobox}

\subsection{Tanpa Pengukuran, Optimasi Hanyalah Dugaan}

Seorang programmer sering merasa yakin bahwa perubahan tertentu pasti
mempercepat program, misalnya mengganti sebuah \asm{MUL} dengan
pergeseran. Keyakinan itu belum tentu salah, tetapi belum terbukti.
Tiga hal membuat dugaan mudah meleset.

\begin{itemize}
    \item \textbf{Hambatan yang berbeda.} Perbaikan hanya berguna bila
          menyerang bagian yang benar-benar paling banyak dieksekusi.
          Mempercepat instruksi yang dijalankan seratus kali tidak
          berarti apa-apa bila ada perulangan yang berjalan sejuta
          kali.
    \item \textbf{Metrik yang bertentangan.} Suatu perubahan dapat
          mempercepat program tetapi memperbesar kode. Bila target
          program adalah ukuran, perubahan tersebut justru merugikan.
    \item \textbf{Lingkungan yang berbeda.} Hasil yang cepat pada
          emulator belum tentu cepat pada perangkat keras
          sesungguhnya, dan sebaliknya.
\end{itemize}

\begin{aktivitasbox}
Ambil satu program yang telah Anda tulis pada bab sebelumnya. Tuliskan
dugaan Anda tentang bagian mana yang paling lambat, lengkap dengan
alasannya. Setelah itu, bongkar program tersebut dengan perintah
\code{U}, lalu bandingkan dugaan Anda dengan kenyataan. Catat pada
jurnal praktikum: berapa kali dugaan Anda benar?
\end{aktivitasbox}
